% SPDX-License-Identifier: Apache-2.0
\section{A Peripheral Written in Software}
\label{sec:x-remote}

A peripheral is usually hardware before it is anything else, and the
program that drives it is written against a datasheet and debugged
against a simulation. \code{Remote} turns that round. It sits on
AXI-Lite where any peripheral would, and answers nothing itself: each
transaction it accepts leaves on a channel as an \code{Ask}, and an
\code{Answer} comes back on another. What is at the far end of those
two channels is not the peripheral's business. In this example it is
a few lines of Rust in the same run; on a board it would be a frame
on a wire and a program on another machine, which is what issue 297
asks for.

So a peripheral can be designed as software, against the bus it will
really sit on, with the rest of the system running as it really runs.
When the software is right, hardware replaces it and nothing else in
the design changes.

What the bus sees is a slow device. One transaction is outstanding at
a time and the channel that carries it is held until the answer
arrives, which is what AXI-Lite allows and what a program on the far
side of a wire needs. A device that never answers must not be able to
stop the bus, so a transaction unanswered for \code{T} cycles is
answered \code{SlvErr} by the peripheral itself. Each carries a tag
and the answer repeats it, which is what makes that safe: an answer
that arrives after the wait ran out carries a tag the peripheral has
finished with, and is dropped rather than given to whatever
transaction is outstanding by then.

The run does four things, in the order a person would try them. It
writes a word and reads it back, so the program holds the state. It
reads a word the program computes rather than stores, the number of
writes it has seen, which is a line of software and a register file
nobody would build. It reads an address the program has no word for,
and the refusal arrives as \code{SlvErr} on the bus. Then it makes
the program slow, twenty cycles a transaction, and every answer still
arrives; then it stops the program answering at all, and the bus is
told the device failed after forty cycles rather than waiting for
ever; then the program comes back and the next read is answered.

The peripheral is lowered, and the build simulates its netlist
against this run under \code{nvc} and Verilator.

\begin{figure}[htbp]
\centering
\input{remote_timing.tex}
\caption{One transaction, end to end: \code{ask} carries it out of the
peripheral, \code{busy} holds while \code{waited} counts, and the
answer comes back and is given to the bus on \code{r}.}
\label{fig:x-remote}
\end{figure}

\lstinputlisting[caption={\code{ex\_remote.rs}. Run with \code{bazel run //lib/examples:ex\_remote}.},label={lst:x-remote}]{lib/examples/ex_remote.rs}

\lstinputlisting[style=ascii,caption={What \code{ex\_remote} printed when the build ran it: the bus and the program, then the lowering.},label={lst:x-remote-out}]{out_remote.txt}
